\setcounter{section}{8}
\Section{Функции многих переменных. Предел и непрерывность}
\Subsection{Метрические пространства и сходимость метрических пространств}

\begin{definition}
Непустое множество \(M\) называется \textit{метрическим пространством}, если
\(\forall x, y \in M\) определено число \(\rho(x, y) \in \R\) такое, что
выполнены аксиомы:
\begin{enumerate}
\item \(\rho(x, y) \geq 0; \; \rho(x, y) = 0 \Longleftrightarrow x = y\).
\item \(\rho(x, y) = \rho(y, x) \quad \forall x, y \in M\).
\item \(\rho(x, y) + \rho(y, z) \geq \rho(x, z) \quad \forall x, y, z \in M\).
\end{enumerate}
\end{definition}

\begin{note}
Расстояние (или метрика) может определяться разными способами. Любое множество
можно превратить в метрическое пространство с тривиальной метрикой:
\[
    \rho(x, y) = \left\{
        \begin{aligned}
            & 0, x = y;\\
            & 1, x \neq y.\\
        \end{aligned}
    \right .\]
\end{note}

\begin{definition}
\textit{Точечное \(n\)-мерное евклидово простраснтво \(\R^n\)}~--- множество
упорядоченных наборов \(x = (\alpha_1, \alpha_2, \dots, \alpha_n), \; \alpha_i
\in \R\).
\end{definition}

\begin{definition}
Если \(x = \alpha_1, \dots, \alpha_n\), \(y = \beta_1, \dots, \beta_n\), то
\textit{расстоянием в~\(\R^n\)} будет величина:
\[
    \rho(x, y) = \sqrt{\sum_{i=1}^n (\alpha_i - \beta_i)^2.}
\]
\end{definition}

\begin{definition}
\(x = y\), если \(\forall i = 1 \dots n \rightarrow \alpha_i = \beta_i\).
\end{definition}

Проверим выполнимость аксиом для введеной нами метрики.

\begin{proof}
Аксиомы 1 и 2 выполняются очевидно.

3. Пусть \(x = (\alpha_1, \dots, \alpha_n)\), \(y = (\beta_1, \dots, \beta_n)\),
\(z \hm = (\gamma_1, \dots, \gamma_n)\). Требуется доказать, что:
\[
    \sqrt{\sum_{i=1}^n (\alpha_i - \beta_i)^2} + \sqrt{\sum_{i=1}^n (\beta_i -
        \gamma_i)^2} \geq \sqrt{\sum_{i=1}^n (\alpha_i - \gamma_i)^2}.
\]

Введём следующие обозначения: \(p_i = \alpha_i - \beta_i\), \(q_i = \beta_i -
\gamma_i\). Тогда \(p_i + q_i = \alpha_i - \gamma_i\), и полученное раннее
неравенство можно переписать следующим образом:
\[
    \sqrt{\sum_{i=1}^n p_i^2} + \sqrt{\sum_{i=1}^n q_i^2} \geq
        \sqrt{\sum_{i=1}^n (p_i + q_i)^2}
\]
поскольку все числа неотрицательны, можно возвести неравенство в квадрат:
\[
    \sum_{i=1}^n p_i^2 + \sum_{i=1}^n q_i^2 + 2\sqrt{\sum_{i=1}^n p_i^2 \cdot
        \sum_{i=1}^n q_i^2} \geq \sum_{i=1}^n (p_i + q_i)^2
\]
раскрыв сумму в правой части, получим:
\[
    \sum_{i=1}^n p_i^2 + \sum_{i=1}^n q_i^2 + 2\sqrt{\sum_{i=1}^n p_i^2 \cdot
        \sum_{i=1}^n q_i^2} \geq \sum_{i=1}^n p_i^2 + \sum_{i=1}^n q_i^2 +
        2\sum_{i=1}^n p_i q_i
\]
\[
    \sqrt{\sum_{i=1}^n p_i^2 \cdot \sum_{i=1}^n q_i^2} \geq \sum_{i=1}^n p_i q_i
\]
последнее неравенство заведомо верно, если:
\[
    \sum_{i=1}^n p_i^2 \cdot \sum_{i=1}^n q_i^2 \geq \left(\sum_{i=1}^n p_i
        q_i\right)^2.
\]

Рассмотрим следующую функцию одной переменной:
\[
    \phi(t) = \sum_{i=1}^n (p_i + q_i)^2 \geq 0 \quad \forall t \in \R
\]
раскроем квадрат:
\[
    \phi(t) = \sum_{i=1}^n p_i^2 + 2\sum_{i=1}^n p_i q_i + \sum_{i=1}^n q_i^2.
\]
поскольку этот квадратный трёхчлен всегда неотрицателен, то:
\[
    D = \left(\sum_{i = 1}^n p_i q_i\right)^2 - \sum_{i=1}^n q_i^2 \cdot
        \sum_{i=1}^n p_i^2 \leq 0.
\]
Откуда следует доказываемое неравенство.
\end{proof}

\begin{note}
В точечных пространствах нет понятия сложения точек и умножения точки на число.
Однако, удобно ввести ноль \(0 = (0, 0, \dots, 0)\). Если \(x = (\alpha_1, \dots,
\alpha_n)\), \(y = (\alpha_1, \dots, \alpha_n)\), то удобно считать, что:
\[
    x \pm y = (\alpha_1 \pm \beta_1, \dots, \alpha_n \pm \alpha_n), \;
        x = \sqrt{\alpha_1^2 + \dots + \alpha_n^2} = \rho(x, 0)
\]
тогда неравенство треугольника перепишется в виде:
\[
    |x - y| + |y - z| \geq |x - z|.
\]

Если \(x - y = p\), \(y - z = q\), то \(x - z = p + q\), и неравенство
перепишется в следующем виде:
\[
    |p| + |q| \geq |p + q|.
\]
\end{note}

\begin{note}
Эти операции удобно делать только в \(\R^n\).
\end{note}

\begin{definition}
\textit{\(\epsilon\)-окрестность точки в метрическом пространстве} (\(\epsilon
\geq 0\)). \[U_\epsilon(x_0) = \{x \in M \colon \rho(x, x_0) < \epsilon\}.\]
\end{definition}

\begin{note}
В случае \(n = 1, 2, 3\) совпадает с обычным понятием.
\end{note}

\begin{note}
В пространстве с тривиальной метрикой при \(\epsilon \leq 1 \; U_\epsilon(x_0)
= \{x_0\}.\). При \(\epsilon > 1 \; U_\epsilon(x_0) = M\).
\end{note}

\begin{definition}
Говорят, что последовательность точек \(x_n\) метрического пространства \(M\)
\textit{сходится к точке \(x_0 \in M\)}, если
\[
    \lim_{n \rightarrow \infty} \rho(x, x_0) = 0, \; \text{т.е.} \; \forall
    \epsilon > 0 \; \exists n_0 \; \forall n \geq n_0 \rightarrow
    \rho(x_n, x_0) < \epsilon.
\]
\end{definition}

В пространстве с тривиальной метрикой \(\ds\lim_{n \rightarrow \infty} x_n = x_0
\Longleftrightarrow \exists n_0 \; \forall n \geq n_0 x_n = x_0\)
(последовательность \textit{стабилизируется}).

\begin{note}
Сохраняются свойства пределов, основанные только на понятии расстояния.
\end{note}

\begin{statement}
Сходящаяся последовательность может иметь только один предел.
\end{statement}

\begin{proof}
Пусть \(\ds \lim_{n \rightarrow \infty} x_n = x_0\) и \(\ds \lim_{n \rightarrow
\infty} x_n = y_0\). Тогда:
\begin{gather}
    \forall \epsilon > 0 \; \exists n_1 \rightarrow \rho(x_n, x_0) <
        \epsilon / 2, \\
    \forall \epsilon > 0 \; \exists n_2 \rightarrow \rho(x_n, y_0) <
        \epsilon / 2.
\end{gather}
Пусть \(n_0 = \max(n_1, n_2)\). Тогда:
\[
    \forall n \geq n_0 \; \rho(x_0, y_0) \leq \rho(x_0, x_n) + \rho(x_n, y_0) <
    \epsilon.
\]
Поскольку \(\epsilon > 0\)~--- любое, то \(\rho(x_0, y_0) = 0 \Rightarrow x_0 =
y_0\).
\end{proof}

\begin{definition}
Пусть \(X\)~--- непустое множество в метрическом пространстве \(M\). Тогда
\textit{диаметром} \(X\) называется величина:
\[
    \diam X = \sup \rho(x, y) \in [0, +\infty] \quad \forall x, y \in X.
\]
\end{definition}

\begin{definition}
Множество \(X\) называется \textit{ограниченным}, если его диаметр конечен.
\end{definition}

\begin{statement}
Множество \(X \subset \R^n\) ограничено \(\Longleftrightarrow \exists C > 0
\colon \forall x \in X \rightarrow |x| \leq C\).
\end{statement}

\begin{proof}
\Boxedrightarrow Пусть множество ограничено, т.е. \(\diam X = d < +\infty\).
Поскольку множество непусто, рассмотрим фиксированное \(x_0 \in X\). Тогда:
\[
    \forall x \in X \rightarrow |x| \leq |x - x_0| + |x_0| \leq d + |x_0|
\]
Тогда \(C = d + |x_0|\).

\Boxedleftarrow Пусть \(\forall x \in X \rightarrow |x| \leq C\). Тогда
\[
    \forall x, y \in X \rightarrow \rho(x, y) = |x - y| \leq |x| + |-y| = |x| +
        |y| \leq 2C \Rightarrow \diam X \leq 2C.
\]
\end{proof}

\begin{definition}
Последовательность \(x_n\) называется \textit{ограниченной}, если множество её
значений ограничено.
\end{definition}

\begin{statement}
Если последовательность \(x_n\) в метрическом пространстве сходится, то она
ограничена.
\end{statement}

\begin{proof}
Пусть \(\ds\lim_{n \Rightarrow \infty} x_n x_0\), т.е.
\begin{align}
    \forall \epsilon > 0 \;\exists n_0 \colon &\forall n \geq n_0 \rightarrow
        \rho(x_n, x_0) < \epsilon / 2 \\
    &\forall m \geq n_0 \rightarrow \rho(x_m, x_0) < \epsilon / 2,
\end{align}
т.е.
\[
    \forall n, m \geq n_0 \rightarrow \rho(x_n, x_m) \leq \rho(x_n, x_0) +
        \rho(x_0, x_m) \leq \epsilon.
\]
Пусть \(\ds C = \max_{n, m \leq n_0} \rho(x_n, x_m) < +\infty\) (конечное число
расстояний). Рассмотрим случай, когда \(n \geq n_0\), а \(m \leq n_0\). Тогда
\[
    \rho(x_n, x_m) \leq \underbrace{\rho(x_n, x_{n_0})}_{< \epsilon} +
        \underbrace{\rho(x_{n_0}, x_m)}_{\leq C} < C + 1.
\]
Получили \(\forall n, m \rightarrow \rho(x_n, x_m) < C + 1\). Следовательно,
последовательность ограничена.
\end{proof}

Рассмотрим сходимость в \(\R^n\). Индекс последовательности обозначим за \(k\).
Пусть \(k\)-й член последовательности это \(x_k = (\alpha_1^k, \dots,
\alpha_n^k)\), \(x_0 = (\alpha_1^0, \dots \alpha_n^0)\)

\begin{statement}
Сходимость в \(\R^n\) равносильна покоординатной сходимости. Т.е.
\[
    x_k \rightarrow x_0 \; \text{в} \; \R^n \Longleftrightarrow \forall i = 1, 2,
        \dots, n \quad \alpha_i^{(k)} \rightarrow \alpha_i^{(0)}.
\]
\end{statement}

\begin{proof}
\Boxedrightarrow Пусть последовательность сходится покоординатно. Докажем, что
имеет место сходимость в привычном смысле.
\[
    \forall \epsilon > 0 \; \exists k_0 \; \forall k \geq k_0 \quad
        \rho(x_k, x_0) < \epsilon \Longleftrightarrow  \forall \epsilon > 0 \;
        \exists k_0 \; \forall k \geq k_0 \quad \sqrt{\sum_{i=1}^n\left(
        x_i^{(k)} - x_i^{(0)}\right)^2}
\]
Ясно, что:
\[
    \forall i = 1 \dots n \quad |x_i^{(k)} - x_i^{(0)}| \leq \sqrt{\sum_{i=1}^n
        \left(x_i^{(k)} - x_i^{(0)}\right)^2} < \epsilon
\]
Это значит, что \(\ds\forall i = 1 \dots n \rightarrow \lim_{k \rightarrow
\infty} x_i^{(k)} = x_i^{(0)}\).

\Boxedleftarrow Пусть \(\ds\forall i = 1 \dots n \rightarrow \lim_{k \rightarrow
\infty} x_i^{(k)} = x_i^{(0)}\). Это значит, что:
\[
    \forall \epsilon > 0 \; \forall i = 1 \dots n \; \exists k_i \colon \forall
        k \geq k_i \rightarrow \left|x_i^{(k)} - x_i^{(0)}\right| <
        \frac{\epsilon}{\sqrt{n}}.
\]
Возьмём \(k_0 = \max k_i, \; i = 1 \dots n\). Тогда:
\[
    \forall \epsilon > 0 \; \exists k_0 \colon \forall k \geq k_0 \rightarrow
        \sqrt{\sum_{i=1}^n \left(x_i^{(k)} - x_i^{(0)}\right)^2} < \sqrt{
        \sum_{i=1}^n\frac{\epsilon^2}{n} \cdot n} = \epsilon
\]
Это значит, что \(\ds\lim_{k \rightarrow \infty} x^{(k)} = x^{(0)}\).
\end{proof}

\begin{definition}
Пусть \(k_m\)~--- последовательность индексов (\(m = 1, 2, \dots\)) такая, что
\(k_m\) возрастает строго. Тогда последовательность \(x_{k_m}\) с индексом \(m\)
называется \textit{подпоследовательностью} последовательности \(x_k\).
\end{definition}

\begin{theorem}[Больцано-Вейерштрасса]
Ограниченная последовательность в \(\R^n\) имеет сходящуюся
подпоследовательность.
\end{theorem}

\begin{proof}
Доказательство для случая \(n = 2\). Пусть имеется последовательность \(x_i =
(\alpha_k, \beta_k)\)~--- ограничена (\(|x_k| \leq C \; \forall k\)). Тогда
также \(\forall k \rightarrow |\alpha_k| \leq C, \; |\beta_k| \leq C\). Применим
к \(\alpha_k\) одномерную теорему Больцано-Вейерштрасса. Это значит, что
\(\exists k_m \colon \alpha_{k_m} \rightarrow \alpha_0\) Последовательность
\(\beta_{k_m}\) ограничена. Из неё выделяем сходящуюся подпоследовательность
\(b_{k_{m_l}} \rightarrow \beta_0\). Тогда \(\alpha_{k_{m_l}} \rightarrow
\alpha_0\) (подпоследовательность сходящейся последовательности имеет тот же
предел). Тогда \((\alpha_{k_{m_l}}, \beta_{k_{m_l}}) \rightarrow
(\alpha_0, \beta_0)\) (покоординатная сходимость равносильна сходимости точек).
Тогда \(x_{k_{m_l}} \rightarrow x_0\).

При \(n > 2\) аналогично (просто получается больше индексов).
\end{proof}

В произвольном метрическом пространстве теорема Больцано-Вейерштрасса неверна.
\begin{proof}
Пусть \(M\)~--- простраснтво с тривиальной метрикой, содержащее счётное число
элементов. Возьмём в качестве последовательности \(x_k\)~--- все элементы \(M\),
как-то занумерованные. Очевидно, это ограниченная последовательность (\(\forall
k, m \; \rho(x_k, x_m) = 1\), т.е. \(\diam \hm = 1\)). Но сходящуюся
подпоследовательность выделить нельзя: любая подпоследовательность такова, что
расстояние между элементами равно единице. Сходящаяся подпоследовательность
стабилизируется, а здесь этого нет.
\end{proof}

\begin{definition}
Последовательность \(x_n\) метрического пространства \(M\) называется
\textit{фундаментальной}, если
\[
    \forall \epsilon > 0 \; \exists n_0 \; \forall n, m \geq n_0 \rightarrow
        \rho(x_n, x_m) < \epsilon.
\]
\end{definition}

\begin{statement}
Сходящаяся последовательность в метрическом пространстве фундаментальна.
\end{statement}

\begin{proof}
Поскольку последовательность сходится, имеем:
\begin{align}
    \forall \epsilon > 0 \;\exists n_0 \colon &\forall n \geq n_0 \rightarrow
        \rho(x_n, x_0) < \epsilon / 2 \\
    &\forall m \geq n_0 \rightarrow \rho(x_m, x_0) < \epsilon / 2,
\end{align}
т.е.
\[
    \forall n, m \geq n_0 \rightarrow \rho(x_n, x_m) \leq \rho(x_n, x_0) +
        \rho(x_0, x_m) < \epsilon.
\]
\end{proof}

Обратное, вообще говоря, неверно. Любое непустое подмножество метрического
пространства можно рассматривать как метрическое пространство с той же метрикой.
\begin{example}
Рассмотрим \(\Q \subset \R, \; \rho(x, y) = |x - y|\). В \(\Q\) можно найти
фундаментальную последовательность, не являющуяся сходящейся.

Рассмотрим \(x_n\)~--- множество десятичных приближений \(\sqrt 2\). \(x_n\)~---
фундаментальна, т.к. сходится в \(\R^n\). В \(\Q\) не является сходящейся. Если
\(\ds\exists \lim_{n \rightarrow \infty} x_n = \alpha\) в \(\Q\), то это предел
и в \(\R\). В силу единственности предела \(\alpha = \sqrt 2 \in \Q\)~---
противоречие.
\end{example}

\begin{definition}
Метрическое пространство, в котором любая фундаментальная последовательность
сходится, называется \textit{полным}.
\end{definition}

\begin{example}
Пространство \(R\)~--- полно, а \(\Q\)~--- неполно (\(\rho(x, y) = |x - y|\)).
\end{example}

\begin{note}
Пространство с тривиальной метрикой полно.
\end{note}

\begin{proof}
Пусть в пространстве с тривиальной метрикой последовательность \(x_n\)
фундаментальна, т.е.:
\[
    \forall \epsilon > 0 \; \exists n_0 \; \forall n, m \geq n_0 \rightarrow
        \rho(x_n, x_m) < \epsilon.
\]
Если взять \(\epsilon = 1\), то \(\forall n, m \geq n_0 \; x_n = x_m\). В
пространстве с тривиальной метрикой последовательность стабилизируется,
следовательно, сходится.
\end{proof}

\begin{theorem}
\(\R^n\)~--- полное пространство.
\end{theorem}

\begin{proof}
Пусть последовательность \(x_k = (\alpha_1^{(k)}, \dots, \alpha_n^{(k)})\)~---
фундаментальна. Это значит, что:
\[
    \forall \epsilon > 0 \; \exists k_0 \colon \forall k, m \geq k_0 \rightarrow
        \rho(x_k, x_m) < \epsilon \Longleftrightarrow \sqrt{\sum_{i=1}^n
        \left(x_i^{(k)} - x_i^{(m)}\right)^2} < \epsilon
\]
\(\forall i = 1 \dots n \; \left|x_i^{(k)} - x_i^{(m)}\right| < \epsilon\), т.е.
\(\ds\forall i = 1 \dots n \rightarrow \exists \lim_{k \rightarrow \infty}
x_i^{(k)} = x_i^{(0)}\). Т.к. покоординатная сходимость равносильна сходимости в
пространстве, то:
\[
    x_k = (x_1^{(k)}, \dots, x_n^{(k)}) \rightarrow (x_1^{(0)}, \dots,
        x_n^{(0)}) \equiv x_0.
\]
\end{proof}


\Subsection{Различные типы множеств в метрических пространствах}

\begin{definition}
Точка \(x_0\) в метрическом пространстве \(M\) называется \textit{внутренней
точкой} множества \(X \subset M\), если
\[
    \exists \epsilon > 0 \colon U_\epsilon(x_0) \subset M.
\]
\end{definition}

\begin{definition}
Множество внутренних точек \(X\) называется \textit{внутренностью} \(X\)
(\(\Int X\)).
\end{definition}

\begin{definition}
Множество \(X \subset M\) называется \textit{открытым}, если \(\Int X = X\), т.е.
любая точка \(X\) является внутренней.
\end{definition}

Ясно, что \(\Int X \subset X\), а совпадает для открытых множеств.

\begin{examples}
В \(\R^1\) открытыми множествами являются открытые промежутки.

В пространстве с тривиальной метрикой для любого множества \(X\) любая точка
является внутренней, т.к.
\[
    \forall x_0 \in X \rightarrow U_\epsilon (x_0) = \{x_0\} \subset X.
\]
\end{examples}

\begin{note}
Пустое множество тоже считается открытым в любом метрическом простраснтве.
\end{note}